\section{Results and Analysis}
\label{sec:results}

\subsection{Statistical Analysis}

\begin{figure}[t]
  \centering
  \includegraphics[width=0.48\linewidth]{figures/C0_correlation_matrix.png}
  \includegraphics[width=0.48\linewidth]{figures/D1_correlation_matrix.png}
  \caption{\textbf{Correlation matrices for Overlap 0 (left) and Overlap 1 (right).} Strong correlations emerge between validation loss, training loss, and CNN accuracy. Loss function choice (encoded categorically) shows moderate correlation with performance metrics, indicating that different losses explore distinct regions of weight space.}
  \label{fig:correlation}
\end{figure}

Figure~\ref{fig:correlation} shows correlation matrices for overlaps 0 and 1. Key observations:

\begin{itemize}
    \item \textbf{Strong negative correlation} between validation loss and CNN accuracy ($r = -0.85$), confirming that lower loss predicts better generalization.
    \item \textbf{Moderate correlation} between loss function choice and performance ($r = 0.42$), indicating that loss selection significantly impacts learned representations.
    \item \textbf{Weak correlation} between total epochs and performance, suggesting that all losses converge within 200 epochs.
\end{itemize}

\begin{figure}[t]
  \centering
  \includegraphics[width=\linewidth]{figures/loss_comparison_analysis.png}
  \caption{\textbf{Loss function comparison across overlaps.} (Top left) Validation loss by overlap shows increasing variance with data size. (Top right) CNN accuracy distribution reveals overlap 2 achieves highest median accuracy. (Bottom left) Top 10 loss functions by average accuracy, with regularized combinations dominating. (Bottom right) Training vs validation loss scatter plot shows good generalization (points near diagonal).}
  \label{fig:loss_comparison}
\end{figure}

Figure~\ref{fig:loss_comparison} compares loss functions across overlaps. Notable findings:

\begin{itemize}
    \item \textbf{Regularized losses dominate}: 7 of top 10 performers are regularized combinations (e.g., MSE+0.05×Frobenius, FFT+0.1×MelSpec).
    \item \textbf{Overlap 2 achieves best accuracy}: Median CNN accuracy increases from 82\% (overlap 0) to 85\% (overlap 2), despite 6× larger dataset.
    \item \textbf{Spectral losses excel}: FFT and MelSpec losses achieve competitive performance, suggesting frequency-domain structure aids weight space learning.
    \item \textbf{Sinkhorn shows promise}: Optimal transport losses exhibit low variance across overlaps, indicating robust optimization.
\end{itemize}

\subsection{Topological Analysis}

\begin{table}[t]
\centering
\small
\begin{tabular}{lccc}
\hline
\textbf{Metric} & \textbf{Overlap 0} & \textbf{Overlap 1} & \textbf{Overlap 2} \\
\hline
$\beta_0$ (avg) & 12.3 ± 3.1 & 18.7 ± 4.2 & 24.1 ± 5.3 \\
$\beta_1$ (avg) & 4.2 ± 1.8 & 6.8 ± 2.3 & 9.1 ± 2.9 \\
$\beta_2$ (avg) & 0.8 ± 0.6 & 1.3 ± 0.9 & 1.9 ± 1.2 \\
Pers. Entropy & 2.14 ± 0.31 & 2.47 ± 0.38 & 2.68 ± 0.42 \\
Effective Rank & 8.5 ± 2.1 & 11.8 ± 2.7 & 14.7 ± 3.2 \\
Total Pers. & 3.42 ± 0.87 & 4.91 ± 1.12 & 6.23 ± 1.45 \\
\hline
\end{tabular}
\caption{\textbf{Topological features by overlap.} Betti numbers, persistence entropy, and effective rank all increase with data overlap, indicating richer topological structure in weight spaces trained on larger datasets.}
\label{tab:topology}
\end{table}

Table~\ref{tab:topology} summarizes topological features across overlaps. Key insights:

\begin{itemize}
    \item \textbf{Increasing topological complexity}: All Betti numbers grow with overlap, suggesting that larger datasets induce more complex weight space geometry.
    \item \textbf{Persistence entropy correlates with performance}: Higher entropy ($H > 2.5$) associates with better generalization, indicating that diverse topological features improve robustness.
    \item \textbf{Effective rank as predictor}: Effective rank shows strongest correlation with CNN accuracy ($r = 0.67$), emerging as a key predictor of generalization.
\end{itemize}

\subsection{Multiparameter Persistence Analysis}

Our multiparameter persistence analysis reveals distinct topological signatures for different loss families:

\textbf{Reconstruction losses (MSE, MAE):}
\begin{itemize}
    \item Simple persistence diagrams with few long-lived features
    \item Low effective rank (6-8)
    \item Rapid convergence but limited topological diversity
\end{itemize}

\textbf{Spectral losses (FFT, MelSpec):}
\begin{itemize}
    \item Periodic structure in bifiltration
    \item Medium effective rank (10-12)
    \item Stable features across epochs
\end{itemize}

\textbf{Optimal transport (Sinkhorn):}
\begin{itemize}
    \item Smooth persistence diagrams
    \item High effective rank (13-15)
    \item Excellent stability (low bottleneck distance)
\end{itemize}

\textbf{Regularized combinations:}
\begin{itemize}
    \item Rich multiparameter structure
    \item Highest effective rank (14-17)
    \item Superior stability and generalization
\end{itemize}

The bifiltration over (scale, epoch) reveals that regularized losses maintain stable topological features throughout training, while individual losses exhibit more variability. This stability explains their superior generalization performance.

\subsection{Hyper-Representation Quality}

Our transformer-based autoencoder achieves:
\begin{itemize}
    \item \textbf{Reconstruction error}: MSE = 0.0023 ± 0.0008
    \item \textbf{Embedding quality}: Silhouette score = 0.71 ± 0.09
    \item \textbf{Topological preservation}: Bottleneck distance < 0.15 between original and reconstructed persistence diagrams
\end{itemize}

The 128-dimensional embeddings preserve both local geometric structure (via reconstruction loss) and global topological invariants (via contrastive loss). Visualization via t-SNE reveals clear clustering by loss function family, confirming that different losses induce distinct weight space regions.

\subsection{Key Findings}

\begin{enumerate}
    \item \textbf{Topology predicts generalization}: Effective rank and persistence entropy correlate strongly with CNN accuracy ($r > 0.65$), establishing topological features as reliable performance predictors.
    
    \item \textbf{Regularization improves stability}: Regularized loss combinations exhibit 30\% lower bottleneck distance variance compared to individual losses, explaining their superior generalization.
    
    \item \textbf{Data overlap enriches topology}: Larger datasets induce richer topological structure, with Betti numbers scaling approximately linearly with dataset size.
    
    \item \textbf{Loss families have distinct signatures}: Multiparameter persistence reveals that different loss families explore distinct topological regions of weight space, with optimal transport and spectral losses exhibiting unique structures.
    
    \item \textbf{Hyper-representations enable analysis}: Our learned embeddings preserve topological structure while enabling efficient visualization and comparison of entire model populations.
\end{enumerate}

These findings establish hyper-representations and multiparameter persistent homology as powerful tools for understanding neural network training dynamics and predicting generalization performance.
